% ==============================================================
\section{問題定義}\label{sec:problem}
% ==============================================================

\subsection{記法}

$\mathcal{I} = \{i_1, i_2, \ldots, i_m\}$ を有限のアイテム集合とする。
$D = [d_1, d_2, \ldots, d_N]$ を時間インデックス $t \in \{1, \ldots, N\}$ で順序付けられたトランザクション列とする。
各トランザクション $d_t \subseteq \mathcal{I}$ はアイテムの集合である。
$k$ 個の補助次元（空間座標、店舗ID等）を考え、各次元は有限定義域 $\mathcal{X}_j$ の値をとる。
各トランザクション $d_t$ は位置 $\mathbf{x}(t) = (x_1(t), \ldots, x_k(t)) \in \mathcal{X} = \prod_j \mathcal{X}_j$ に関連付けられる。

\subsection{サポート曲面}

\begin{definition}[多次元サポート曲面]\label{def:support_surface}
パターン $P$、時間窓 $W_0$、空間窓 $W_j$（$j = 1, \ldots, k$）に対し、$P$ の\emph{サポート曲面} $S_P : \mathbb{Z}^{k+1} \to \mathbb{Z}_{\geq 0}$ を
\begin{equation}\label{eq:support_surface}
S_P(t, \mathbf{v}) = \bigl|\{d_r \in D \mid t \leq r < t + W_0,\; P \subseteq d_r,\; \mathbf{x}(r) \in \mathrm{Box}(\mathbf{v}, \mathbf{W})\}\bigr|
\end{equation}
と定義する。$k = 0$ のとき標準的な1次元サポート時系列に帰着する。
\end{definition}

\subsection{外部イベントの時空間拡張}

1次元設定のイベント $e = (\textit{id}, \textit{name}, t_s, t_e)$ を時空間に拡張する。

\begin{definition}[時空間外部イベント]\label{def:st_event}
\emph{時空間外部イベント} $e = (\textit{id}, \textit{name}, t_s, t_e, \mathcal{V}_e)$ は、
識別子、名称、時間区間 $[t_s, t_e]$、および\emph{空間スコープ}
$\mathcal{V}_e \subseteq \mathcal{X}$（イベントが影響を及ぼす空間領域）から成る。
$\mathcal{V}_e = \mathcal{X}$（全域）のとき、1次元設定のイベント定義に帰着する。
イベント集合を $\mathcal{E} = \{e_1, \ldots, e_K\}$ とする。
\end{definition}

\begin{example}
スーパーマーケットチェーンにおいて、「渋谷・新宿エリア限定 夏季キャンペーン」は
$e = (\text{C01}, \text{夏季}, 150, 250, \{\text{渋谷}, \text{新宿}\})$ と表現される。
このイベントは店舗次元で$\{$渋谷, 新宿$\}$、時間区間$[150, 250]$にのみ影響する。
\end{example}

\subsection{時空間変化点}

1次元設定では変化点は閾値交差の時刻 $\tau$ として定義される。
多次元設定では、変化点に空間的位置が付随する。

\begin{definition}[位置条件付きサポート時系列]\label{def:local_series}
空間位置 $\mathbf{v} \in \mathcal{X}$ における\emph{位置条件付きサポート時系列}は
$s_P^{\mathbf{v}}(t) = S_P(t, \mathbf{v})$
であり、サポート曲面を空間位置 $\mathbf{v}$ で切り出した1次元時系列である。
\end{definition}

\begin{definition}[時空間変化点]\label{def:st_changepoint}
位置 $\mathbf{v}$ において、$s_P^{\mathbf{v}}$ の閾値 $\theta$ に関する上昇交差時刻 $\tau$ を
\emph{時空間変化点} $(\tau, \mathbf{v})$ と呼ぶ。
すなわち、$s_P^{\mathbf{v}}(\tau) \geq \theta$ かつ $(\tau = 0$ または $s_P^{\mathbf{v}}(\tau-1) < \theta)$ が成立する。
下降交差も同様に定義する。
パターン $P$ の時空間変化点の集合を $\mathcal{C}_P = \{(\tau_i, \mathbf{v}_i, \textit{dir}_i)\}$ と表す。
\end{definition}

\begin{remark}[変化領域]
同一パターンが近い時刻に複数の空間位置で変化点を持つとき、
これらは空間的に一貫した変化を表す。
変化点集合 $\{(\tau_i, \mathbf{v}_i)\}$ のうち時間的に近い（$|\tau_i - \tau_j| \leq \delta_t$）かつ
空間的に隣接する点の連結成分を\emph{変化領域}と呼ぶ。
変化領域は「いつ・どこで変化が起きたか」を空間解像度で記述する。
\end{remark}

\subsection{時空間帰属スコア}

\begin{definition}[時空間帰属スコア]\label{def:st_score}
パターン $P$ の時空間変化点 $(\tau, \mathbf{v})$ とイベント $e$ に対し、
\emph{時空間帰属スコア}を
\begin{equation}\label{eq:st_score}
A(P, \tau, \mathbf{v}, e) = \mathrm{prox}_t(\tau, e) \cdot \mathrm{prox}_s(\mathbf{v}, e) \cdot \mathrm{mag}(\tau, \mathbf{v})
\end{equation}
と定義する。3つの構成要素は以下の通りである。

\textbf{時間的近接度}:
\begin{equation}\label{eq:prox_t}
\mathrm{prox}_t(\tau, e) = \exp\!\left(-\frac{\min(|\tau - t_s|, |\tau - t_e|)}{\sigma_t}\right)
\end{equation}

\textbf{空間的近接度}:
\begin{equation}\label{eq:prox_s}
\mathrm{prox}_s(\mathbf{v}, e) = \begin{cases}
1 & \text{if } \mathbf{v} \in \mathcal{V}_e \\
\exp\!\left(-\frac{d_s(\mathbf{v}, \mathcal{V}_e)}{\sigma_s}\right) & \text{otherwise}
\end{cases}
\end{equation}
ここで $d_s(\mathbf{v}, \mathcal{V}_e) = \min_{\mathbf{u} \in \mathcal{V}_e} \|\mathbf{v} - \mathbf{u}\|_1$ はイベント空間スコープまでの$L_1$距離である。

\textbf{変化量}:
\begin{equation}\label{eq:st_mag}
\mathrm{mag}(\tau, \mathbf{v}) = |\bar{s}^{+}(\tau, \mathbf{v}) - \bar{s}^{-}(\tau, \mathbf{v})|
\end{equation}
$\bar{s}^{\pm}$ は変化点前後の局所平均サポートである（1次元設定と同じ定義を各位置で適用）。
\end{definition}

\begin{remark}[1次元設定への帰着]
$k = 0$（空間次元なし）かつ $\mathcal{V}_e = \mathcal{X}$ のとき、
$\mathrm{prox}_s \equiv 1$ となり、式~\eqref{eq:st_score}は
1次元パイプラインのスコア $A = \mathrm{prox}_t \cdot \mathrm{mag}$ に帰着する。
\end{remark}

パターン $P$ のイベント $e$ に対する観測帰属スコアは、
全変化点にわたる総和とする：
\begin{equation}\label{eq:obs_score}
A_{\mathrm{obs}}(P, e) = \sum_{(\tau, \mathbf{v}) \in \mathcal{C}_P} A(P, \tau, \mathbf{v}, e).
\end{equation}

\subsection{時空間イベント帰属問題}

\begin{problem}[時空間イベント帰属]\label{prob:st_attribution}
\textbf{入力}：位置付きトランザクション列 $D$、
窓サイズ $(W_0, \mathbf{W})$、閾値 $\theta$、最大パターン長 $L_{\max}$、
時空間イベント集合 $\mathcal{E}$、有意水準 $\alpha$。

\textbf{出力}：$s_P^{\mathbf{v}}$ の閾値交差変化点がイベント $e$ の時間窓および空間スコープと
統計的に有意に関連する全ての組 $(P, e)$ を、
偽発見率（FDR）を水準 $\alpha$ 以下に制御しながら同定する。
\end{problem}

本問題の主要な課題は以下の4点である：
\begin{enumerate}
\item \textbf{仮説空間の爆発}：$|\mathcal{F}| \times |\mathcal{E}|$ 個のパターン--イベント対に加え、
各対が複数の空間位置で独立に帰属しうるため、多重検定の負荷が1次元設定より大きい。

\item \textbf{空間自己相関}：近接する空間位置のサポート時系列は正に相関するため、
変化点が空間的にクラスター化し、独立性仮定に基づく検定は過度に楽観的になる。

\item \textbf{副次パターン問題}：1次元設定と同様、アイテムを共有するパターン間でサポートが連動変動する。

\item \textbf{イベントの空間的部分影響}：イベントが一部の空間領域にのみ影響する場合、
全域で集約したサポートでは信号が希釈される（空間的混同）。
\end{enumerate}
